\documentclass[../../main.tex]{subfiles}

\begin{document}

\section{积拓扑}\label{section:ch02-积拓扑}

本章余下部分将继续讨论如何在集合上给出拓扑的种种方法。

在此之前，我们曾为两个拓扑空间的积 $X \times Y$ 定义了拓扑，本节将这个定义推广到任意笛卡儿积。

为此我们考虑笛卡儿积
\begin{align*}
X_{1} \times \cdots \times X_{n} \quad \text{和} \quad X_{1} \times X_{2} \times \cdots,
\end{align*}
其中 $X_{i}$ 都是拓扑空间。定义笛卡儿积的拓扑有两种方式。第一种方法是分别以形如 $U_{1} \times \cdots \times U_{n}$ 和 $U_{1} \times U_{2} \times \cdots$ 的集合作为上述相应笛卡儿积的拓扑的基元素，其中对于每一个 $i$，$U_{i}$ 表示 $X_{i}$ 中的一个开集。用这种方式为笛卡儿积定义的拓扑，我们将称之为箱拓扑。

第二种方法是将第\cref{section:ch02-积拓扑XY}节中给出的用子基定义拓扑的方式予以推广。也即以所有形如 $\pi_{i}^{-1}(U_{i})$ 的集合构成子基，这里 $i$ 是任意给定的指标，$U_{i}$ 是 $X_{i}$ 中的一个开集。这种子基生成的拓扑我们称之为积拓扑。

两种拓扑有什么不同？考虑第二种拓扑的典型基元素 $B$。它是子基元素 $\pi_{i}^{-1}(U_{i})$ 的有限交，比如说 $i = i_{1}, \cdots, i_{k}$。则点 $x$ 属于 $B$ 当且仅当对于 $i = i_{1}, \cdots, i_{k}$ 时 $\pi_{i}(x) \in U_{i}$，而当 $i$ 为其余值时，对于 $\pi_{i}(x)$ 则没有任何限制。

这意味着两种拓扑对于有限笛卡儿积是相同的，对于无限笛卡儿积是不同的。为何我们更偏爱第二种拓扑呢？这是我们本节将要说明的问题。

首先我们要引进更为广义的笛卡儿的概念。在此之前，我们仅就指标集为 $\{1, \cdots, n\}$ 或者为 $\mathbb{Z}_{+}$ 的情形定义了加标族的笛卡儿积。以下我们考虑指标集为任意集合的情形。

\begin{definition}[$J$-串与笛卡儿积]\label{definition:ch02-J串}
设 $J$ 是一个指标集。对于给定的集合 $X$，$X$ 的元素的 $J$-串（$J$-tuple）定义为一个映射 $x: J \rightarrow X$。若 $\alpha$ 为 $J$ 的一个元素，我们用 $x_{\alpha}$ 表示 $x$ 在 $\alpha$ 处的值，而再不用 $x(\alpha)$ 表示这个值，并且将它称为 $x$ 的第 $\alpha$ 个坐标（coordinate）。我们将用记号
\begin{align*}
(x_{\alpha})_{\alpha \in J}
\end{align*}
表示函数 $x$ 本身。这个“串记法”使我们将 $J$-串对于指标族的依赖表示得更明确些。用 $X^{J}$ 表示 $X$ 中元素的 $J$-串的全体。

设 $\{A_{\alpha}\}_{\alpha \in J}$ 是一个加标集族，$X = \bigcup_{\alpha \in J} A_{\alpha}$。加标集族 $\{A_{\alpha}\}_{\alpha \in J}$ 的笛卡儿积（Cartesian product），定义为使得对于每一个 $\alpha \in J$ 有 $x_{\alpha} \in A_{\alpha}$ 的 $X$ 的元素的所有 $J$-串 $(x_{\alpha})_{\alpha \in J}$ 的集合，用
\begin{align*}
\prod_{\alpha \in J} A_{\alpha}
\end{align*}
表示。也就是说，它是所有这样的函数
\begin{align*}
x: J \longrightarrow \bigcup_{\alpha \in J} A_{\alpha}
\end{align*}
的集合，这些函数要求满足条件：对于每一个 $\alpha \in J$ 有 $x(\alpha) \in A_{\alpha}$。

当指标集无需强调时，有时我们也用 $\prod A_{\alpha}$ 表示上述笛卡儿积，其元素则记为 $(x_{\alpha})$。

当所有的 $A_{\alpha}$ 都等于同一个集合 $X$ 时，笛卡儿积 $\prod_{\alpha \in J} A_{\alpha}$ 恰为所有 $X$ 的元素的 $J$-串的集合 $X^{J}$。对于 $X^{J}$ 的元素，我们有时使用“串记法”表示，有时使用函数表示，依方便而定。
\end{definition}

\begin{definition}[箱拓扑与积拓扑]\label{definition:ch02-箱拓扑积拓扑}
设 $\{X_{\alpha}\}_{\alpha \in J}$ 是拓扑空间的一个加标族。积空间
\begin{align*}
\prod_{\alpha \in J} X_{\alpha}
\end{align*}
上的某一个拓扑的基取为所有形如
\begin{align*}
\prod_{\alpha \in J} U_{\alpha}
\end{align*}
的集合的族，其中对于每一个 $\alpha \in J$，$U_{\alpha}$ 在 $X_{\alpha}$ 中是开的。由这个基生成的拓扑叫做箱拓扑（box topology）。

因为 $\prod X_{\alpha}$ 本身是一个基元素，所以 $\{\prod U_{\alpha}\}$ 满足基定义中的第一个条件。又因为任意两个基元素的交是另一个基元素
\begin{align*}
\left(\prod_{\alpha \in J} U_{\alpha}\right) \cap \left(\prod_{\alpha \in J} V_{\alpha}\right) = \prod_{\alpha \in J} (U_{\alpha} \cap V_{\alpha}),
\end{align*}
所以它满足基定义中的第二个条件。

下面的方法是用子基定义拓扑的推广。设函数
\begin{align*}
\pi_{\beta}: \prod_{\alpha \in J} X_{\alpha} \longrightarrow X_{\beta}
\end{align*}
将笛卡儿积空间的每一个元素对应其第 $\beta$ 个坐标，即
\begin{align*}
\pi_{\beta}((x_{\alpha})_{\alpha \in J}) = x_{\beta},
\end{align*}
我们把 $\pi_{\beta}$ 称为关于指标 $\beta$ 的投射（projection mapping）。

令 $\mathcal{S}_{\beta}$ 表示族
\begin{align*}
\mathcal{S}_{\beta} = \{\pi_{\beta}^{-1}(U_{\beta}) \mid U_{\beta} \text{ 在 } X_{\beta} \text{ 中是开的}\},
\end{align*}
令 $\mathcal{S}$ 为所有族 $\mathcal{S}_{\beta}$ 的并
\begin{align*}
\mathcal{S} = \bigcup_{\beta \in J} \mathcal{S}_{\beta}.
\end{align*}
子基 $\mathcal{S}$ 生成的拓扑称为积拓扑（product topology）。给定了这个拓扑的 $\prod_{\alpha \in J} X_{\alpha}$ 称为积空间（product space）。
\end{definition}

为比较积拓扑与箱拓扑，只要考虑由 $\mathcal{S}$ 生成的基 $\mathcal{B}$。族 $\mathcal{B}$ 是由 $\mathcal{S}$ 中元素的所有有限交组成的。但是，如果我们取属于同一个 $\mathcal{S}_{\beta}$ 中的元素进行交，那并不能得到什么新东西，因为
\begin{align*}
\pi_{\beta}^{-1}(U_{\beta}) \cap \pi_{\beta}^{-1}(V_{\beta}) = \pi_{\beta}^{-1}(U_{\beta} \cap V_{\beta});
\end{align*}
$\mathcal{S}_{\beta}$ 中的两个或有限多个元素的交，还是 $\mathcal{S}_{\beta}$ 的一个元素。只有取不属于同一个 $\mathcal{S}_{\beta}$ 的元素进行交，才能得到一些新东西。因此，基 $\mathcal{B}$ 的典型元素可以这样描述：令 $\beta_{1}, \cdots, \beta_{n}$ 为指标集 $J$ 中不同指标的一个有限集，$U_{\beta_{i}}$ 为 $X_{\beta_{i}}$ 中的一个开集（$i = 1, \cdots, n$）。则
\begin{align*}
B = \pi_{\beta_{1}}^{-1}(U_{\beta_{1}}) \cap \pi_{\beta_{2}}^{-1}(U_{\beta_{2}}) \cap \cdots \cap \pi_{\beta_{n}}^{-1}(U_{\beta_{n}})
\end{align*}
是 $\mathcal{B}$ 的一个典型的元素。

因此，点 $x = (x_{\alpha})$ 在 $B$ 中当且仅当 $x$ 的第 $\beta_{1}$ 个坐标在 $U_{\beta_{1}}$ 中，第 $\beta_{2}$ 个坐标在 $U_{\beta_{2}}$ 中等等。如果指标 $\alpha$ 不是 $\beta_{1}, \cdots, \beta_{n}$ 中的一个，那么 $x$ 的第 $\alpha$ 个坐标就没有任何限制，于是可以将 $B$ 写成积的形式
\begin{align*}
B = \prod_{\alpha \in J} U_{\alpha},
\end{align*}
其中，当 $\alpha \neq \beta_{1}, \cdots, \beta_{n}$ 时，$U_{\alpha}$ 表示空间 $X_{\alpha}$。

综上所述可得下面的定理：

\begin{theorem}[箱拓扑与积拓扑的比较]\label{theorem:ch02-箱积比较}
$\prod X_{\alpha}$ 的箱拓扑以形如 $\prod U_{\alpha}$ 的集合作为基元素，其中，对于每一个 $\alpha$，$U_{\alpha}$ 在 $X_{\alpha}$ 中是开的。$\prod X_{\alpha}$ 的积拓扑以形如 $\prod U_{\alpha}$ 的集合作为基元素，其中 $U_{\alpha}$ 在 $X_{\alpha}$ 中是开的，并且除去有限多个 $\alpha$ 外，对于每一个 $\alpha$ 都有 $U_{\alpha} = X_{\alpha}$。

有两件事是很清楚的。第一，对于有限积 $\prod X_{\alpha}$，两种拓扑是一样的。第二，一般说来，箱拓扑细于积拓扑。
\end{theorem}

那么为什么我们常常用积拓扑而不用箱拓扑呢？答案要从我们对拓扑学的研究中去找。我们发现，如果使用积拓扑，那么关于有限积空间的一些重要定理对于任意积空间也成立；而使用箱拓扑则不然。因此，积拓扑在数学中更为重要，而箱拓扑就不那么重要。但是它可以用来构造反例。所以：当我们讨论积空间 $\prod X_{\alpha}$ 的时候，如果不特别申明，总是假定所给的就是积拓扑。

在有关积空间 $X \times Y$ 的已经证明过的定理之中，有一些对于积空间 $\prod X_{\alpha}$ 不论采用箱拓扑或积拓扑都成立。下面列出一些这样的定理，大多数相关的证明留作习题。

\begin{theorem}[由基构造箱拓扑与积拓扑]\label{theorem:ch02-基构造积}
设每一个空间 $X_{\alpha}$ 的拓扑由基 $\mathcal{B}_{\alpha}$ 给出。则形如
\begin{align*}
\prod_{\alpha \in J} B_{\alpha}
\end{align*}
的集族是 $\prod_{\alpha \in J} X_{\alpha}$ 的箱拓扑的一个基，其中对于每一个 $\alpha$，$B_{\alpha} \in \mathcal{B}_{\alpha}$。

对于如上形式的集族，如果仅对有限多个指标 $\alpha$ 要求 $B_{\alpha} \in \mathcal{B}_{\alpha}$，而对余下的指标有 $B_{\alpha} = X_{\alpha}$，则这个集族便是 $\prod_{\alpha \in J} X_{\alpha}$ 的积拓扑的一个基。
\end{theorem}

\begin{example}\label{example:ch02-Rn基}
考虑 $n$ 维欧氏空间 $\mathbb{R}^{n}$。$\mathbb{R}$ 中所有开区间组成 $\mathbb{R}$ 的一个基，因此所有形如
\begin{align*}
(a_{1}, b_{1}) \times (a_{2}, b_{2}) \times \cdots \times (a_{n}, b_{n})
\end{align*}
的积组成了 $\mathbb{R}^{n}$ 的一个拓扑基。因为 $\mathbb{R}^{n}$ 是有限积，箱拓扑与积拓扑是一样的。所以在讨论 $\mathbb{R}^{n}$ 的时候，除非特别申明，都采用上面给出的这种拓扑。
\end{example}

\begin{theorem}[积空间的子空间]\label{theorem:ch02-积子空间}
设对于每一个 $\alpha \in J$，$A_{\alpha}$ 是 $X_{\alpha}$ 的一个子空间，则当两者都用箱拓扑或者两者都用积拓扑时，$\prod A_{\alpha}$ 是 $\prod X_{\alpha}$ 的一个子空间。
\end{theorem}

\begin{theorem}[积空间的 Hausdorff 性质]\label{theorem:ch02-积Hausdorff}
若每一个空间 $X_{\alpha}$ 都是 Hausdorff 的，则无论是箱拓扑还是积拓扑，$\prod X_{\alpha}$ 都是 Hausdorff 的。
\end{theorem}

\begin{theorem}[积空间的闭包]\label{theorem:ch02-积闭包}
设 $\{X_{\alpha}\}$ 是一个加标空间族，对每一个 $\alpha$ 有 $A_{\alpha} \subset X_{\alpha}$。若对 $\prod X_{\alpha}$ 赋予积拓扑或者赋予箱拓扑，则有
\begin{align*}
\prod \overline{A_{\alpha}} = \overline{\prod A_{\alpha}}.
\end{align*}
\end{theorem}

\begin{proof}
设 $x = (x_{\alpha})$ 为 $\prod \overline{A_{\alpha}}$ 的一个点。我们来证明 $x \in \overline{\prod A_{\alpha}}$。设 $U = \prod U_{\alpha}$ 为箱拓扑或是积拓扑空间中含有 $x$ 的一个基元素。由于 $x_{\alpha} \in \overline{A_{\alpha}}$，对于每一个 $\alpha$，可以选取点 $y_{\alpha} \in U_{\alpha} \cap A_{\alpha}$。于是 $y = (y_{\alpha})$ 既属于 $U$ 又属于 $\prod A_{\alpha}$。由于 $U$ 的任意性，$x$ 属于 $\prod A_{\alpha}$ 的闭包。

反之，设对于两个拓扑中的任何一个，$x = (x_{\alpha})$ 属于 $\prod A_{\alpha}$ 的闭包。我们来证明对于任何一个指标 $\beta$，有 $x_{\beta} \in \overline{A_{\beta}}$。设 $V_{\beta}$ 为包含 $x_{\beta}$ 的 $X_{\beta}$ 中的开集。由于对于两个拓扑中的无论哪一个来说，$\pi_{\beta}^{-1}(V_{\beta})$ 是 $\prod X_{\alpha}$ 中的开集，所以它含有 $\prod A_{\alpha}$ 中的点 $y = (y_{\alpha})$。于是 $y_{\beta}$ 属于 $V_{\beta} \cap A_{\beta}$。从而 $x_{\beta} \in \overline{A_{\beta}}$。
\end{proof}

到目前为止，我们仍然没有理由只喜欢积拓扑，而不要箱拓扑；但是当我们试图推广前面关于映到积空间的映射的连续性的那个定理时，这种区别就出现了。当赋予 $\prod X_{\alpha}$ 箱拓扑时，下面定理不成立。

\begin{theorem}[到积空间的映射（积拓扑）]\label{theorem:ch02-到积空间映射}
映射 $f: A \rightarrow \prod_{\alpha \in J} X_{\alpha}$ 定义为
\begin{align*}
f(a) = (f_{\alpha}(a))_{\alpha \in J},
\end{align*}
其中对每一个 $\alpha$，$f_{\alpha}: A \rightarrow X_{\alpha}$。设 $\prod X_{\alpha}$ 具有积拓扑，则 $f$ 连续当且仅当每一个函数 $f_{\alpha}$ 连续。
\end{theorem}

\begin{proof}
设 $\pi_{\beta}$ 是积空间到其第 $\beta$ 个坐标空间上的投射。$\pi_{\beta}$ 是连续的，这是因为如果 $U_{\beta}$ 是 $X_{\beta}$ 的一个开集，集合 $\pi_{\beta}^{-1}(U_{\beta})$ 就是 $\prod X_{\alpha}$ 的积拓扑的一个子基元素。假定 $f: A \rightarrow \prod X_{\alpha}$ 连续，则 $f_{\beta} = \pi_{\beta} \circ f$ 是两个连续函数的复合，所以 $f_{\beta}$ 连续。

反之，设每一个坐标函数 $f_{\alpha}$ 都连续。要证明 $f$ 连续，只要证明每一个子基元素在 $f$ 下的原像是 $A$ 的开集就够了。在我们定义连续函数的时候已经提到过这一点。对于 $\prod X_{\alpha}$ 的积拓扑，其典型子基元素是 $\pi_{\beta}^{-1}(U_{\beta})$，其中 $\beta$ 是某一个指标，$U_{\beta}$ 是 $X_{\beta}$ 的开集。因为 $f_{\beta} = \pi_{\beta} \circ f$，所以 $f^{-1}(\pi_{\beta}^{-1}(U_{\beta})) = f_{\beta}^{-1}(U_{\beta})$。又因为 $f_{\beta}$ 连续，因此这个集合是 $A$ 的一个开集。
\end{proof}

为什么这个定理对于箱拓扑不成立呢？或许最令人信服的办法就是考察一个例子。

\begin{example}\label{example:ch02-箱拓扑不连续}
考虑 $\mathbb{R}$ 的可数无限积 $\mathbb{R}^{\omega}$。
\begin{align*}
\mathbb{R}^{\omega} = \prod_{n \in \mathbb{Z}_{+}} X_{n},
\end{align*}
对于每一个 $n \in \mathbb{Z}_{+}$，$X_{n} = \mathbb{R}$。函数 $f: \mathbb{R} \rightarrow \mathbb{R}^{\omega}$ 定义为
\begin{align*}
f(t) = (t, t, \cdots),
\end{align*}
其第 $n$ 个坐标函数是 $f_{n}(t) = t$。每一个坐标函数 $f_{n}: \mathbb{R} \rightarrow \mathbb{R}$ 都连续，所以当赋予 $\mathbb{R}^{\omega}$ 积拓扑时，$f$ 是连续的。而当 $\mathbb{R}^{\omega}$ 给的是箱拓扑时，$f$ 不连续。例如，取箱拓扑的基元素
\begin{align*}
B = (-1, 1) \times \left(-\frac{1}{2}, \frac{1}{2}\right) \times \left(-\frac{1}{3}, \frac{1}{3}\right) \times \cdots,
\end{align*}
我们断言：$f^{-1}(B)$ 不是 $\mathbb{R}$ 中开集。事实上，如果 $f^{-1}(B)$ 是 $\mathbb{R}$ 中开集，它必定包含 0 旁边的某一个区间 $(-\delta, \delta)$，也就是说 $f((-\delta, \delta)) \subset B$。将 $\pi_{n}$ 作用于这个包含关系的两边，可见对于所有 $n$，
\begin{align*}
f_{n}((-\delta, \delta)) = (-\delta, \delta) \subset \left(-\frac{1}{n}, \frac{1}{n}\right),
\end{align*}
这是一个矛盾。
\end{example}

\subsection*{习题}

\begin{exercise}\label{exercise:ch02-积拓扑-1}
证明\cref{theorem:ch02-基构造积}。
\end{exercise}

\begin{exercise}\label{exercise:ch02-积拓扑-2}
证明\cref{theorem:ch02-积子空间}。
\end{exercise}

\begin{exercise}\label{exercise:ch02-积拓扑-3}
证明\cref{theorem:ch02-积Hausdorff}。
\end{exercise}

\begin{exercise}\label{exercise:ch02-积拓扑-4}
证明：$(X_{1} \times \cdots \times X_{n-1}) \times X_{n}$ 与 $X_{1} \times \cdots \times X_{n}$ 同胚。
\end{exercise}

\begin{exercise}\label{exercise:ch02-积拓扑-5}
\cref{theorem:ch02-到积空间映射}的陈述中，有一个蕴涵关系对于箱拓扑也成立，是哪一个蕴涵关系？
\end{exercise}

\begin{exercise}\label{exercise:ch02-积拓扑-6}
设 $x_{1}, x_{2}, \cdots$ 是积空间 $\prod X_{\alpha}$ 的点的一个序列。证明：这个序列收敛到点 $x$ 当且仅当对于每一个 $\alpha$，序列 $\pi_{\alpha}(x_{1}), \pi_{\alpha}(x_{2}), \cdots$ 收敛到 $\pi_{\alpha}(x)$。若用箱拓扑代替积拓扑相应结论还成立吗？
\end{exercise}

\begin{exercise}\label{exercise:ch02-积拓扑-7}
设 $\mathbb{R}^{\infty}$ 是 $\mathbb{R}^{\omega}$ 中所有“终端为 0”的序列（即，使得仅有有限多个 $i$，$x_{i} \neq 0$）的所有序列 $(x_{1}, x_{2}, \cdots)$ 的子集。在箱拓扑与积拓扑下，$\mathbb{R}^{\infty}$ 在 $\mathbb{R}^{\omega}$ 中的闭包是什么？验证你的结论。
\end{exercise}

\begin{exercise}\label{exercise:ch02-积拓扑-8}
给定实数序列 $(a_{1}, a_{2}, \cdots)$ 和 $(b_{1}, b_{2}, \cdots)$，其中对于所有的 $i$，$a_{i} > 0$，用
\begin{align*}
h((x_{1}, x_{2}, \cdots)) = (a_{1}x_{1} + b_{1}, a_{2}x_{2} + b_{2}, \cdots)
\end{align*}
定义一个函数 $h: \mathbb{R}^{\omega} \rightarrow \mathbb{R}^{\omega}$。证明：若赋予 $\mathbb{R}^{\omega}$ 积拓扑，则 $h$ 是 $\mathbb{R}^{\omega}$ 的一个自同胚。若赋予 $\mathbb{R}^{\omega}$ 箱拓扑，结论会怎样？
\end{exercise}

\begin{exercise}\label{exercise:ch02-积拓扑-9}
证明选择公理等价于以下条件：对于每一由非空集合组成的加标集族 $\{A_{\alpha}\}_{\alpha \in J}$，$J \neq \varnothing$，其笛卡儿积
\begin{align*}
\prod_{\alpha \in J} A_{\alpha}
\end{align*}
非空.
\end{exercise}

\begin{exercise}\label{exercise:ch02-积拓扑-10}
设 $A$ 是一个集合，$\{X_{\alpha}\}_{\alpha \in J}$ 是空间的一个加标族，$\{f_{\alpha}\}_{\alpha \in J}$ 是函数 $f_{\alpha}: A \rightarrow X_{\alpha}$ 的一个加标族。

(a) 证明：$A$ 有唯一的一个使得每一个 $f_{\alpha}$ 都连续的最粗拓扑 $\mathcal{T}$。

(b) 设
\begin{align*}
\mathcal{S}_{\beta} = \{f_{\beta}^{-1}(U_{\beta}) \mid U_{\beta} \text{ 在 } X_{\beta} \text{ 中是开的}\},
\end{align*}
并且 $\mathcal{S} = \bigcup \mathcal{S}_{\beta}$。证明：$\mathcal{S}$ 是 $\mathcal{T}$ 的一个子基。

(c) 证明：映射 $g: Y \rightarrow A$ 关于 $\mathcal{T}$ 连续当且仅当每一个映射 $f_{\alpha} \circ g$ 都连续。

(d) 用
\begin{align*}
f(a) = (f_{\alpha}(a))_{\alpha \in J}
\end{align*}
定义映射 $f: A \rightarrow \prod X_{\alpha}$。设 $Z$ 表示积空间 $\prod X_{\alpha}$ 的子空间 $f(A)$。证明 $\mathcal{T}$ 中每一个元素在 $f$ 下的像是 $Z$ 中的一个开集。
\end{exercise}

\end{document}